\chapter{Conclusion and Future directions}
In this work, we have investigated critical phenomena in axisymmetric vacuum spacetimes using Brill wave initial data with twist. Through the construction and evolution of one-parameter families of initial data, we identified a threshold separating dispersion from black hole formation. For the family of pure twisting Brill waves constructed here, the critical parameter has been tuned to a precision of three decimal places. 
While previous studies typically achieve higher precision (of order $\sim 6$ decimal places), the present level of tuning corresponds to an intermediate stage in the approach to the critical threshold. Nevertheless, it already allows us to observe a scaling relation between the maximum of the Kretschmann scalar and the distance from the critical parameter. An estimate of the critical exponent can also be obtained from this scaling behaviour. However, given the current level of tuning, these results should be regarded as preliminary, and further refinement is required before drawing more definitive conclusions.

Additionally, near-critical evolutions exhibit echo-like features in the curvature, suggestive of underlying self-similar dynamics. While the current level of tuning and resolution does not allow for a definitive identification of discrete self-similarity, these results are consistent with expectations from critical collapse. 

We have also discussed the numerical challenges encountered in this work and possible strategies to mitigate them. In particular, the use of a Cartesian-based apparent horizon finder has already improved the exploration of the parameter space and is expected to facilitate further tuning toward the critical regime.
The substantial computational cost ( $\approx$4000 core-hours per simulation) further highlights the difficulty of accessing the critical regime in less symmetric spacetimes.

\section*{Future Directions}
An immediate direction for future work is to achieve higher precision in tuning the critical parameter. This would enable more accurate determination of critical exponents and a clearer assessment of universality. Exploring multiple families of initial data will be essential to establish whether the observed behaviour is truly universal. Until now, we have only tried pure twisting Brill wave i.e. Seed function $q(\rho,z) =0 $ and after trying different families in that, we will move on to adding the influence of the seed function along with twist variables. Then an even stronger test would be to take the other class of axisymmetric vacuum initial data which are Teukolsky waves, modify them to include twist and conduct similar tests in that data. We expect that to be easier with the insights that we have gained with this experiment.

A much deeper investigation of self-similarity is also required. Higher resolution simulations may allow for the extraction of echoing structure and scaling periods, providing stronger evidence for the nature of the critical solution.

Understanding the role of twist remains an open problem. In particular, it is important to determine whether twist gives rise to new classes of critical solutions or modifies existing ones. This could provide insight into the broader structure of the solution space in axisymmetric vacuum spacetimes.

From a numerical perspective, improving gauge choices and developing more robust horizon-finding techniques will be crucial. We have already implemented a new gauge choice that appears to facilitate reaching the horizon more efficiently; however, it is not discussed here, as it requires further testing and validation. 
Additionally, the Cartesian-based apparent horizon finder currently requires a substantial amount of manual intervention. Developing automated or more robust horizon-finding methods will therefore be an important direction for future work.

